\subsection{Integral curves}

In this section, X denotes a manifold of class \(C^{r}\) and \(\xi\) is a vector field of class \(C^{s}\) on X, with r, s in \(N_{K}\) , \(s \leqslant r - 1\) .

9.1.1. Let I be an open subset of K, and let \(f: \mathbf{I} \to \mathbf{X}\) a map of class \(\mathbf{C}^k\) ( \(k \in \mathbf{N}_{\mathbb{K}}\) , \(k \leqslant r\) ) from I into X. For \(t \in \mathbf{I}\) , one denotes by \(f'(t)\) the vector \(\mathrm{T}_t(f)(1) \in \mathrm{T}_{f(t)}(\mathbf{X})\) ; this vector is sometimes called the velocity of \(f\) at time \(t\) . When X is a Banach space, \(f'(t)\) is the derivative of \(f\) at \(t\) .

One says that \(f\) is an integral curve of \(\xi\) if one has

\[
f ^ {\prime} (t) = \xi (f (t))\tag{1}
\]

for all \(t \in I\) . Such a curve is of class \(C^{s+1}\) .

If \(x\) is a point of X, one calls integral curve of \(\xi\) with origin \(x\) an integral curve \(f: I \to X\) of \(\xi\) such that \(0 \in I\) and \(f(0) = x\) ; for all \(x \in X\) , there exists an integral curve of \(\xi\) with origin \(x\) , defined in a suitable open neighborhood of 0 in K.

If \(f: \mathbf{I} \to \mathbf{X}\) is an integral curve of \(\xi\) and if \(a \in \mathbf{K}\) , the map \(t \mapsto f(t - a)\) of \(\mathbf{I} + a\) to \(\mathbf{X}\) is an integral curve of \(\xi\) .

Let \(f_{1}\colon\mathbf{I}_{1}\to\mathbf{X}\) and \(f_{2}\colon\mathbf{I}_{2}\to\mathbf{X}\) two integral curves of \(\xi\) , and \(t\in\mathbf{I}_{1}\cap\mathbf{I}_{2}\) . If \(f_{1}(t)=f_{2}(t)\) the maps \(f_{1}\) and \(f_{2}\) coincide in a neighborhood of \(t\) .

9.1.2 ("Flows"). Let W be an open subset of X × K. For every x ∈ X, denote by W \(_{x}\) the set of t ∈ K such that (x, t) ∈ W. We call an integral flow of ξ with domain W a map f of class C \(^{k}\) (k ≤ r) from W to X such that, for every x ∈ X, the map f \(_{x}\) : W \(_{x}\) → X defined by f \(_{x}\) (t) = f(x, t) is an integral curve of ξ, and that whenever (x, 0) ∈ W, one has f(x, 0) = x.

For every point \(x \in X\) , there exists an integral flow of \(\xi\) of class \(\mathbf{C}^s\) whose domain is a neighborhood of \((x, 0)\) in \(X \times K\) . Two integral flows of \(\xi\) defined in a neighborhood of \((x, 0)\) coincide on a neighborhood of \((x, 0)\) .

Let \(x\in X\) and let \(f\) be an integral flow of \(\xi\) whose domain W is a neighborhood of \((x,0)\) . There exists a neighborhood V of \((x,0,0)\) in \(X \times K \times K\) such that, if \((x,t_{1},t_{2})\in\mathrm{V}\) , one has

\[
(x, t _ {1}) \in \mathbf {W}, \quad (f (x, t _ {1}), t _ {2}) \in \mathbf {W}, \quad (x, t _ {1} + t _ {2}) \in \mathbf {W}
\]

and

\[
f (f (x, t _ {1}), t _ {2}) = f (x, t _ {1} + t _ {2}).\tag{2}
\]

If X is paracompact, or if X is separated and K = R or C, there exists an integral flow of \(\xi\) whose domain is a neighborhood of \(X \times \{0\}\) .

9.1.3 ("The real case"). We assume that K = R and that X is separated. We call an integral arc of \(\xi\) any integral curve of \(\xi\) whose domain is an open interval of R. Two integral arcs \(f_{1}\) and \(f_{2}\) of \(\xi\) , defined in intervals \(I_{1}\) and \(I_{2}\) , and coinciding at a point of \(I_{1} \cap I_{2}\) , coincide in \(I_{1} \cap I_{2}\) .

For every point \(x \in X\) , there exists an integral arc with origin \(x\) and one and only one \(f_x: I_x \to X\) such that, for every integral arc \(f: I \to X\) with origin \(x\) , one has \(I \subset I_x\) and \(f = f_x|I\) . The arc \(f_x\) is called the maximal integral arc of \(\xi\) with origin \(x\) , and its domain \(I_x\) is sometimes called the lifetime interval of the point \(x\) for the field \(\xi\) .

If \(I_x = ]-a_-,a_+[\) (and if \(a_+\) is finite, the restriction of \(f_x\) to the interval \([0,a_+[\) is a proper map from \([0,a_+[\) into X. In particular, \(f_x(t)\) has no limit when \(t\) tends to \(a_+\) .

9.1.4. With the notations and hypotheses of 9.1.3, the set \(\Omega\) of the pairs \((x,t)\in\mathbf{X}\times\mathbf{R}\) such that \(t\in\mathbf{I}_{x}\) is open in \(\mathbf{X}\times\mathbf{R}\) and the map \(f\colon\Omega\to\mathbf{X}\) defined by \(f(x,t)=f_{x}(t)\) is an integral flow of \(\xi\) with domain \(\Omega\) .

The functions \(\alpha_{+}\) and \(\alpha_{-}: X \to ]0,+\infty]\) defined by \(I_{x} = ]-\alpha_{-}(x),\alpha_{+}(x)[\) are lower semicontinuous. For \((x, t) \in \Omega\) and \(y = f(x, t)\) , one has \(\alpha_{+}(y) = \alpha_{+}(x) - t\) and \(\alpha_{-}(y) = \alpha_{-}(x) + t\) .

If \(t_1 \in \mathbf{I}_x\) and if \(t_2 \in \mathbf{I}_{f(x, t_1)}\) , one has \(t_1 + t_2 \in \mathbf{I}_x\) and

\[
f (f (x, t _ {1}), t _ {2}) = f (x, t _ {1} + t _ {2}).\tag{3}
\]

If the function \(\alpha_{+}\) (resp. \(\alpha_{-}\) ) is bounded below on X by a constant \(>0\) , it is constant and equal to \(+\infty\) .

One has \(\alpha_{+} = \alpha_{-} = +\infty\) (in other words, one has \(\Omega = \mathbf{X} \times \mathbf{R}\) ) in each of the following cases:

\begin{enumerate}
\def\labelenumi{(\roman{enumi})}
\item
  the field \(\xi\) has compact support (for example if X is compact);
\item
  there exists a transitive group of automorphisms of X preserving \(\xi\) ;
\item
  there exists on X a Riemannian structure for which X is complete and \(\xi\) bounded.
\end{enumerate}

If \(\Omega = X \times R\) , the map \(f: X \times R \to X\) is a right action law of class \(C^s\) of the group manifold \(R\) on the manifold \(X\) , cf.~5.12.5.

9.1.5 ("The complex case"). Suppose that K = C and that X is separated. For every \(a \in ]0, +\infty]\) , denote by \(D_a\) the open disk with center 0 and radius a in C. For every \(x \in X\) , there exists one and only one number \(\rho(x) \in ]0, +\infty]\) and one and only one integral curve \(f_x: D_{\rho(x)} \to X\) with origin \(x\) such that, for every \(a \in ]0, +\infty]\) and every integral curve \(f: \mathrm{D}_a \to \mathrm{X}\) with origin \(x\) , one has \(a \leqslant \rho(x)\) and \(f = f_x | \mathrm{D}_a\) .

For every \(\lambda\in\mathbf{C}\) let \(\alpha_{\lambda}(x)\) be the upper bound of the lifetime interval of \(x\) for the vector field \(\lambda\xi\) on the real manifold deduced from X by restriction of scalars. One has \(\rho(x)=\inf_{|\lambda|=1}\alpha_{\lambda}(x)\) .

The set \(\Delta\) of the pairs \((x,t)\in X\times\mathbf{C}\) such that \(t\in D_{\rho(x)}\) is open in \(X\times\mathbf{C}\) and the map \(f:\Delta\to X\) defined by \(f(x,t)=f_x(t)\) is an integral flow of \(\xi\) . The function \(\rho:X\to ]0,+\infty]\) is lower semicontinuous. For \((x,t)\in\Delta\) and \(y=f(x,t)\) , one has \(\rho(y)\geqslant\rho(x)-|t|\) . If \(\rho\) is bounded below on \(X\) by a constant \(>0\) , one has \(\rho = +\infty\) and \(\Delta = X \times C\) .

When \(\Delta = X \times C\) , the map \(f: X \times C \to X\) is a right action law of the group manifold \(C\) on the manifold \(X\) .

9.1.6. Let \(\varphi: X \to Y\) be a morphism of manifolds and \(\eta\) a vector field on Y such that \(\xi\) is \(\varphi\)-related to \(\eta\) (cf. 8.2.6). If \(f: I \to X\) is an integral curve of \(\xi\) , the map \(\varphi \circ f: I \to Y\) is an integral curve of \(\eta\) . If \(K = R\) and if X and Y are separated, for every \(x \in X\) , the lifetime interval of \(x\) for \(\xi\) is contained in the lifetime interval of \(\varphi(x)\) for \(\eta\) . If moreover \(\varphi\) is proper, these intervals are equal.

9.1.7 ("Time-dependent equations"). Let W be an open subset of X × K and let Ξ: W → T(X) be a lifting of class C \(^{s}\) of pr \(_{1}\) : W → X, that is, a map of class C \(^{s}\) from W into T(X) such that Ξ(x, t) belongs to T \(_{x}\) (X) for all (x, t) ∈ W. One calls an integral curve of Ξ a map f of class C \(^{k}\) (with k ∈ N \(_{K}\) and k ≤ r), from an open subset I of K into X such that, for all t ∈ I, one has

\[
(f (t), t) \in \mathrm{W},\qquad f ^ {\prime} (t) = \Xi (f (t), t).
\]

Let \(\eta\) be the vector field of class \(C^s\) on W defined by \(\eta(x, t) = (\Xi(x, t), (t, 1))\). For a map \(f\) from an open subset I of K into X to be an integral curve of \(\Xi\), it is necessary and sufficient that \(t \mapsto (f(t), t)\) be an integral curve of \(\eta\) in the sense of 9.1.1.

Let \(g: \mathbf{I} \to \mathbf{W}\) be an integral curve of \(\eta\) and \(t_0\) a point of \(\mathbf{I}\) such that \(\operatorname{pr}_2(g(t_0)) = t_0\). The set of elements \(t \in \mathbf{I}\) such that \(\operatorname{pr}_2(g(t)) = t\) is a neighborhood of \(t_0\) that coincides with \(\mathbf{I}\) if \(\mathbf{I}\) is connected.

9.1.8 ("Parameters and initial conditions"). Let Z be a manifold of class \(C^{r}\) , let W be an open subset of X × K × Z , and let \(\Xi\) be a lifting of class \(C^{s}\) of \(pr_{1}: W \to X\). If \(z \in Z\) , one denotes by \(W_{z}\) the set of \((x, t) \in X \times K\) such that \((x, t, z) \in W\) , and one denotes \(\Xi_{z}\) as the map \((x, t) \mapsto \Xi(x, t, z)\) from \(W_{z}\) into T(X).

Let \((x_0, t_0, z_0)\) be a point of W. There then exists an open neighborhood \(X_1 \times I_1 \times Z_1\) of \((x_0, t_0, z_0)\) in W and a map of class \(C^s\)

\[
f \colon X _ {1} \times I _ {1} \times I _ {1} \times Z _ {1} \rightarrow X
\]

such that, for any \((x_{1}, t_{1}, z_{1}) \in \mathbf{X}_{1} \times \mathbf{I}_{1} \times \mathbf{Z}_{1}\) , the map \(t \mapsto f(x_{1}, t_{1}, t, z_{1})\) is an integral curve of \(\Xi_{z_{1}}\) (in the sense of 9.1.7) satisfying the relation \(f(x_{1}, t_{1}, t_{1}, z_{1}) = x_{1}\) .

Two maps

\[
f _ {1}: \mathrm{X} _ {1} \times \mathrm{I} _ {1} \times \mathrm{I} _ {1} \times \mathrm{Z} _ {1} \rightarrow \mathrm{X} \quad \text {and} \quad f _ {2}: \mathrm{X} _ {2} \times \mathrm{I} _ {2} \times \mathrm{I} _ {2} \times \mathrm{Z} _ {2} \rightarrow \mathrm{X}
\]

satisfying the above conditions coincide in a neighborhood of \((x_{0}, t_{0}, t_{0}, z_{0})\) .

